\begin{theorem}[Endpoint conductors in wild prime degree]
\label{mod:cases-wildendpoints}
Let $K/F$ be a wildly ramified cyclic extension of prime degree
$p=[K:F]$, with positive lower break $t$, residue degree one, and a chosen
integral uniformizer $\varpi_K$ generating $\mathcal O_K$ over
$\mathcal O_F$.  Put $\varpi_F=N_{K/F}(\varpi_K)$ and
$\psi_K=\psi_F\circ\operatorname{Tr}_{K/F}$.  If
$n=n_F(\psi_F)$, the exact trace-conductor formula used in both endpoint
calculations is
\[
 n_K(\psi_K)=pn+(p-1)(t+1).
\]
Every product below ranges over the complete norm-character group
$S(K/F)$, including its identity element.

Suppose first that $m_F(\chi_F)=0$.  The identity norm character has
conductor zero, while every nonidentity $\mu\in S(K/F)$ has conductor
$t+1$.  Direct substitution in the finite local-constant sum, with exact
uniformizer-power denominators, gives
\[
 \Delta_F(\mu\chi_F,\psi_F)
 =\chi_F(\varpi_F)^{m_F(\mu)+n}
   \Delta_F(\mu,\psi_F).
\]
The identity norm-character factor is proved literally equal to $1$, and
the identity twist is the original $\Delta_F(\chi_F,\psi_F)$.  Summing
the exact conductors over the full group gives
\[
 n+(p-1)(t+1+n)=pn+(p-1)(t+1),
\]
and hence the complete conductor-zero product identity
\[
 \Delta_K(\chi_F\circ N_{K/F},\psi_K)
 \prod_{\mu\in S(K/F)}\Delta_F(\mu,\psi_F)
 =
 \prod_{\mu\in S(K/F)}\Delta_F(\mu\chi_F,\psi_F).
\]
No stationary class or Lamprecht formula is used in this branch.

Now suppose that $m_F(\chi_F)=1$.  Lean chooses an admissible denominator
and normalizes its residual additive character to the canonical
Frobenius-invariant character.  If $f$ is the minimal polynomial of
$\varpi_K$, it sets
\[
 \gamma_K=\gamma_F\frac{f'(\varpi_K)}{\varpi_K^{p-1}}.
\]
The different calculation gives
$v_K(f'(\varpi_K))=(p-1)(t+1)$, and the trace-dual calculation retains the
exact sign
\[
 \operatorname{Tr}_{K/F}
   \left(\frac{\varpi_K^{p-1}}{f'(\varpi_K)}\right)=1.
\]
Together with the trace-ideal bound, this identifies the residual additive
characters.  The conductor-one residue formula is applied on both fields
with its leading minus sign and inverse multiplicative character.  The
Frobenius change of variables then yields
\[
 \frac{\Delta_K(\chi_F\circ N_{K/F},\psi_K)}
      {\Delta_F(\chi_F,\psi_F)}
 =\chi_F\left(\frac{N_{K/F}(\gamma_K)}{\gamma_F}\right).
\]

Let $\lambda\in k_F^\times$ be the normalized ramification displacement.
The generic positive-break construction in
Theorem~\ref{mod:ramification-normbelowbreak} fixes the exact source and
target coordinates of the critical graded norm and proves
\[
 P(Z)=Z^p-\lambda^{p-1}Z,
 \qquad
 \operatorname{range}(P)
   =\lambda^p\ker\bigl(\operatorname{Tr}_{k_F/\mathbf F_p}\bigr).
\]
It also records exactness of the ramification map followed by the critical
norm and fixes the cokernel in the target-over-range direction:
\[
 \#\left(
   (U_F^t/U_F^{t+1})/
   \operatorname{range}(\operatorname{gr}N_{K/F})
 \right)=p.
\]
The endpoint declarations specialize that shared construction at
$t=s+1$; there is no second endpoint-specific proof of the polynomial or
its linear coefficient.

Let $\tau$ generate the nontrivial norm-character line.  Its conductor is
$t+1>1$, so Lean chooses a representative $b_\tau$ of its genuine
stationary numerator class and puts
$c_\tau=\gamma_\tau/b_\tau$, where
$\gamma_\tau=\varpi_F^t\gamma_F$.  Triviality of $\tau$ on the actual
critical norm image, the exact polynomial, and the trace-pairing
annihilator give
\[
 \left(
   \overline{\frac{c_\tau}{\varpi_F^t\gamma_F}}
 \right)^{p-1}
 =\lambda^{p(p-1)}.
\]
For
$D=f'(\varpi_K)/\varpi_K^{(p-1)(t+1)}$, the conjugate-difference product
and Wilson's theorem give
\[
 \bar D=(-1)^p\lambda^{p-1},
 \qquad
 \overline{N_{K/F}(D)}=(-1)^p\lambda^{p(p-1)}.
\]
Consequently Lean proves the endpoint discriminant congruence in the exact
quotient direction
\[
 \frac{N_{K/F}(\gamma_K)/\gamma_F}
      {(-1)^p c_\tau^{p-1}}
 \in U_F^1.
\]

For every nonzero cyclic index $j$, stationary-class power calculus proves
that $j b_\tau$ represents the class of $\tau^j$.  Stable twisting is
applied only to these conductor-$t+1$ norm characters and gives the factor
$c_\tau/j$.  Their complete product has residue
$(-1)^p c_\tau^{p-1}$, including when $p=2$.  The identity norm character
is retained and evaluated separately.  Combining these facts proves the
complete conductor-one product identity
\[
 \Delta_K(\chi_F\circ N_{K/F},\psi_K)
 \prod_{\mu\in S(K/F)}\Delta_F(\mu,\psi_F)
 =
 \prod_{\mu\in S(K/F)}\Delta_F(\mu\chi_F,\psi_F).
\]

\LeanFile{LanglandsFirstMainLemma/Cases/WildEndpoints.lean}
\LeanFile{LanglandsFirstMainLemma/Cases/WildEndpoints/Common.lean}
\LeanFile{LanglandsFirstMainLemma/Cases/WildEndpoints/Zero.lean}
\LeanFile{LanglandsFirstMainLemma/Cases/WildEndpoints/One/ResidualGauss.lean}
\LeanFile{LanglandsFirstMainLemma/Cases/WildEndpoints/One/CyclicWilson.lean}
\LeanFile{LanglandsFirstMainLemma/Cases/WildEndpoints/One/DerivativeDiscriminant.lean}
\LeanFile{LanglandsFirstMainLemma/Cases/WildEndpoints/One/IntrinsicCritical.lean}
\LeanFile{LanglandsFirstMainLemma/Cases/WildEndpoints/One/StationaryReciprocity.lean}
\LeanFile{LanglandsFirstMainLemma/Cases/WildEndpoints/One/Final.lean}
\LeanFile{LanglandsFirstMainLemma/FiniteField/FrobeniusTrace.lean}
\lean{LanglandsFirstMainLemma.firstMain_wildEndpoint_zero}
\lean{LanglandsFirstMainLemma.firstMain_wildEndpoint_one}
\lean{LanglandsFirstMainLemma.endpoint_discriminant_congruence}
\lean{LanglandsFirstMainLemma.wildEndpointIntrinsicCriticalExactAndCokernel}
\lean{LanglandsFirstMainLemma.wildEndpointIntrinsicCriticalPolynomialValue_exact}
\lean{LanglandsFirstMainLemma.wildEndpointIntrinsicCriticalPolynomialValue_range_eq_traceKernel}
\leanok
\SourceText{Subsection sec:wild-endpoints and Lemma lem:endpoint-discriminant-congruence.}
\uses{mod:ramification-normfiltration,
mod:ramification-normbelowbreak,
mod:ramification-normcharacters,
mod:ramification-pullbackconductors,
mod:ramification-traceideals,
mod:delta-residueformula,
mod:lamprecht-stabletwist}
\FormalizationContract
The two principal Lean theorems are literal \texttt{deltaFinite} product
identities with proof-bearing exact conductor data and internally fixed
admissible denominators.  They do not assert an identity for an arbitrary
parameter of type \texttt{LocalConstantFunction}; that requires a separate
bridge from arbitrary continuous characters to exact conductor data.
Every product uses the complete finite norm-character group.  Its identity
term is never erased: it is proved equal to $1$ on the norm-character side
and to the original $\chi_F$ factor on the twisted side.  The critical
polynomial declarations are endpoint wrappers around the generic
positive-break API in module~\ref{mod:ramification-normbelowbreak}.  The
critical cokernel is the target modulo the actual norm range, not the
reverse quotient.  The residue Gauss sum keeps the manuscript's leading
minus sign and inverse-character convention, and the discriminant congruence keeps
the displayed quotient direction.  Stationary numerator classes are
constructed only for nontrivial norm characters of conductor $t+1>1$;
neither endpoint character of conductor zero or one is assigned a
stationary parameter or evaluated by a nonexistent higher-conductor
Lamprecht layer.
\end{theorem}
